\documentclass[11pt,a4paper]{article}
\usepackage[utf8]{inputenc}
\usepackage[T1]{fontenc} % Vital para la correcta codificación de tildes y caracteres especiales (|)
\usepackage[spanish,es-tabla]{babel}
\usepackage{amsmath, amssymb}
\usepackage{geometry}
\geometry{top=2.5cm, bottom=2.5cm, left=2.5cm, right=2.5cm}
\usepackage{fancyhdr}
\usepackage{circuitikz}
\usepackage{siunitx} % Estándar internacional para unidades de ingeniería

% Configuración de siunitx para el contexto de ingeniería
\sisetup{
    output-decimal-marker = {.},
    inter-unit-product = \ensuremath{{}\cdot{}}
}

\pagestyle{fancy}
\fancyhf{}
\rhead{IMT-221 \textbar{} Evaluación Diagnóstica}
\lhead{Ingeniería Mecatrónica \textbar{} UCB Sede Tarija}
\cfoot{\thepage}

\begin{document}

\begin{center}
    \textbf{\Large UNIVERSIDAD CATÓLICA BOLIVIANA ``SAN PABLO'' - SEDE TARIJA}\\[0.3cm]
    \textbf{\large Carrera de Ingeniería Mecatrónica \textbar{} Circuitos Electrónicos II (IMT-221)}\\[0.5cm]
    \textbf{\Large EVALUACIÓN DIAGNÓSTICA - GESTIÓN II-2026}
\end{center}

\vspace{0.5cm}
\noindent \textbf{Estudiante:} \rule{8cm}{0.4pt} \hfill \textbf{Fecha:} 03 de agosto de 2026

% Se reemplazó el salto '\\' problemático por un \vspace seguro.
\vspace{0.5cm}
\noindent \textbf{Instrucciones:} Resuelva los siguientes ejercicios de forma manuscrita. Esta prueba no tiene ponderación sumativa en la nota continua, pero es un requisito obligatorio para medir las competencias previas del grupo.

\vspace{0.5cm}
\noindent \textbf{Pregunta 1: Álgebra de Números Complejos (25\%)} \\
Dado el número complejo en forma rectangular $Z = -6 + j8$:
\begin{enumerate}
    \item Conviértalo a su forma polar ($R \angle \theta$) y a su forma exponencial ($R e^{j\theta}$).
    \item Calcule su inverso complejo $1/Z$ en forma rectangular.
\end{enumerate}

\vspace{0.5cm}
\noindent \textbf{Pregunta 2: Análisis de Nodos en CC (25\%)} \\
Para el siguiente circuito de corriente continua, determine la tensión en el nodo $V_A$ aplicando la Ley de Corrientes de Kirchhoff (LCK):

\begin{center}
\begin{circuitikz}[american, scale=1.2, transform shape]
    % Malla principal exterior
    % Se utiliza 'l' para la etiqueta (nombre) y 'a' para la anotación (valor)
    \draw (0,0) to[V, l=$V_s$, a=\SI{24}{\volt}] (0,2) 
                to[R, l=$R_1$, a=\SI{4}{\ohm}] (3,2) 
                to[R, l=$R_2$, a=\SI{12}{\ohm}] (6,2)
                to[I, l=$I_o$, a=\SI{2}{\ampere}] (6,0) 
                -- (0,0);
                
    % Rama central con polos de conexión automáticos (*-*)
    \draw (3,2) to[R, l=$R_3$, a=\SI{6}{\ohm}, *-*] (3,0);
    
    % Referencia a tierra (GND) aislada para evitar conflictos de trazo
    \draw (0,0) node[ground]{};
    
    % Etiqueta del nodo VA declarada de forma independiente
    \node[above=0.1cm] at (3,2) {$V_A$};
\end{circuitikz}
\end{center}

\vspace{0.5cm}
\noindent \textbf{Pregunta 3: Relación Voltaje-Corriente en Capacitores (25\%)} \\
Un capacitor de $C = \SI{4.7}{\micro\farad}$ se encuentra inicialmente descargado ($v_c(0) = \SI{0}{\volt}$). Si se le inyecta un pulso de corriente constante de $i_c(t) = \SI{10}{\milli\ampere}$ durante un intervalo de $t = 0$ a $t = \SI{5}{\milli\second}$:
\begin{enumerate}
    \item Calcule la tensión final en los bornes del capacitor $v_c(\SI{5}{\milli\second})$.
    \item Determine la energía almacenada en el campo eléctrico en ese instante ($W = \frac{1}{2} C v^2$).
\end{enumerate}

\vspace{0.5cm}
\noindent \textbf{Pregunta 4: Parámetros de Ondas Senoidales (25\%)} \\
Una señal de tensión de corriente alterna capturada en un osciloscopio está descrita por:
$$v(t) = 170 \cos(314.16 t - 45^\circ) \, [\unit{\volt}]$$
\begin{enumerate}
    \item Indique el valor pico de tensión ($V_p$) y la frecuencia de la señal en Hertz ($f$).
    \item Calcule la tensión Eficaz o RMS ($V_{RMS}$).
\end{enumerate}

\end{document}